\section{Descending diagrams and a constructive unknotting bound}
\label{sec:descending}

Choose a regular basepoint and a traversal direction on a one-component
diagram. The diagram is \emph{descending} if the first visit to each
crossing is on its overstrand. We shall turn this condition on the
$2c$-letter visit word into a certificate of polygonal unknotting in
three-space.

\begin{lemma}[Changing the heights of a lift]
Two rational polygonal lifts of the same parametrized regular diagram,
with the same overpass at each crossing, are related by rational
triangle moves and subdivisions in three-space.
\end{lemma}
\begin{proof}
Subdivide the two projected circuits simultaneously at all height
breakpoints and at both occurrences of every crossing. A vertex at a
crossing now has a distinct vertex of the other branch above or below it.
If a vertex $A$ with neighbors $U,V$ is moved vertically to $A'$ without
sweeping past or meeting that other branch, two triangle moves implement
the displacement. Insert $A'$ on $U A$ using the triangle $UAA'$, and
delete $A$ from $A'AV$ using the triangle $A'AV$. Skip zero displacements.
Both triangles are nondegenerate: the neighboring projections are
different from that of $A$. Each projects onto one incident projected
halfedge. Its open projected halfedge contains no other branch, by the
subdivision. Over the moved vertex the triangle occupies precisely the
interval of heights being swept; over the other endpoint it occupies
just the unchanged endpoint. The two incident halfedges have disjoint
interiors. These observations verify the clean-triangle tests, including
when a neighboring vertex is another crossing occurrence.

At each crossing first move the upper vertex to the larger of its
initial and target heights and the lower to the smaller of its two
heights. Then move them to their target heights. The
upper vertex can be moved first, since its target is strictly higher
than the target lower height and than the temporarily lowered vertex.
The lower vertex can then be moved. Every intermediate comparison is
strict. Move ordinary vertices once to their targets. This gives at most
four triangle moves per common-subdivision vertex.
\end{proof}

\begin{theorem}[A descending diagram bounds a polygonally unknotted lift]
Every rational polygonal lift of a descending diagram admits a finite
rational triangle-move sequence to a triangle. Both a lift and such
a sequence are primitive recursively constructible.
\end{theorem}
\begin{proof}
Take the basepoint $b$ in the interior of an edge, and choose points
$b^-,b^+$ immediately before and after it in traversal order.
The short interval from $b^-$ to $b^+$ contains no crossing, bends or
other passages of the diagram. Its endpoints can be chosen by taking
midpoints to the nearest recorded event on either side along that edge.
Subdivide at these two points, all old bends and all crossing occurrences.
Read the resulting planar list from $b$ back to $b$ as
$p_0,p_1,\ldots,p_m$, with $p_0=p_m=b$, $p_1=b^+$ and $p_{m-1}=b^-$.
Choose strictly decreasing heights $z_i=m-i$ and put
$A_i=(p_i,z_i)$ and $B_i=(b,z_i)$. First form the temporary polygon
consisting of the chain $A_0\cdots A_m$ and the vertical return
$A_m A_0$. The chain is simple by strict descent; the return meets it
only at the two ends since $b$ has no other preimage. At every crossing
the earlier passage is higher.

To obtain a lift with the exact regular projection, replace the local
arc $A_{m-1}A_m A_0 A_1$ by the straight segment $A_{m-1}A_1$.
First delete $A_m$ across the triangle $A_{m-1}A_mA_0$, then delete
$A_0$ across $A_{m-1}A_0A_1$. These triangles project into
$[b^-,b^+]$, away from the rest of the diagram except at its prescribed
ends. The first triangle lies over $[b^-,b]$ and the remaining piece
$A_0A_1$ lies over $[b,b^+]$. The second triangle is nondegenerate
because $A_0$ is higher than both its other vertices. The new straight
segment traces the short basepoint interval once. Thus the new polygon
really projects to the original diagram with its descending crossing
bits. Reversing these two moves returns to the temporary polygon.

Subdivide its vertical return at the $B_i$. Between heights $z_i$ and
$z_{i+1}$ take the two triangles
\[
 (B_i,A_i,A_{i+1}),\qquad (B_i,A_{i+1},B_{i+1}),
\]
omitting the repeated-vertex end triangles. They form an embedded strip.
For completeness, at the fraction $0<s<1$ down the slab the diagonal
has planar coordinate $b+s(p_{i+1}-b)$. The second triangle runs in
that slice from $b$ to this point; the first continues from it by the
vector $(1-s)(p_i-b)$ to the chain. Two nonzero directions
$p_i-b$ and $p_{i+1}-b$ cannot be negative multiples, since then the
interior of the edge $p_i p_{i+1}$ would pass through $b$. Consequently
the two successive segments in the slice meet only at their common end.
If one vector is zero we are in an end slab, containing just one
triangle. Different open slabs have disjoint heights and consecutive
ones meet along the common horizontal rung $B_i A_i$.

Peel this strip from the top. In slab zero delete $B_0=A_0$ across the
second triangle, leaving the rung $B_1A_1$. At a subsequent step replace
$B_i A_i A_{i+1}$ by $B_i A_{i+1}$ across the first triangle, then
$B_{i+1}B_i A_{i+1}$ by $B_{i+1}A_{i+1}$ across the second. The embedded
strip and level separation give clean moves at each step. Stop with
the last triangle $B_{m-1}A_{m-1}A_m$. Thus at most $2m+2$ moves after
subdivision suffice for the chosen lift. The previous lemma transfers
the certificate to any other lift over the same parametrized projection.
\end{proof}

\begin{theorem}[Warping degree and the crossing-change bound]
For an oriented diagram $D$ with $c>0$ crossings, let $d(D)$ be the
minimum, over its $2c$ basepoint gaps, of the number of first visits
which are underpasses. Write $-D$ for the opposite traversal. Then
\[
                d(D)+d(-D)\le c-1,
\]
with equality exactly when the cyclic visits alternate over and under.
One can change at most $\lfloor(c-1)/2\rfloor$ crossings of $D$ to
obtain a diagram whose rational polygonal lifts are triangle-equivalent
to a triangle. At $c=0$ the corresponding cost is zero.
\end{theorem}
\begin{proof}
Let $f(j)$ count first-under crossings starting in gap $j$. For the same
gap the opposite traversal sees the other visit first at every crossing,
so its count is $c-f(j)$. Advancing the basepoint just past one visit
affects a single crossing. When the passed visit is over, it ceases to
be first and the under visit becomes first, so the count rises by one.
Passing an under visit lowers it by one. Thus
\[
 d(D)+d(-D)=c-\bigl(\max_j f(j)-\min_j f(j)\bigr).
\]
The range is at least one. It equals one exactly when the unit steps
alternate: two successive equal steps give range at least two, and
alternating steps visit just two adjacent integers. This proves the
claimed equality criterion. Choose a minimizing gap in whichever
direction has smaller minimum, and switch all its first-under
crossings. This changes at most $\lfloor(c-1)/2\rfloor$ bits and yields
a descending traversal. The previous theorem gives its PL certificate;
for $c=0$ every traversal is already descending. The chosen traversal
simply constructs the crossing bits and height order and may be used
regardless of which orientation is retained on the input knot.
\end{proof}

\begin{corollary}[Unlinking a diagram]
Any $k$-component diagram with $c$ crossings can be changed at at most
$\lfloor c/2\rfloor$ crossings to one admitting an explicit reduction by
spatial triangle moves and subdivisions to $k$ triangles in disjoint
horizontal slabs.
\end{corollary}
\begin{proof}
Order the components; choose a direction and a regular basepoint on each.
Request earlier components to pass over later ones at mixed crossings,
and request the first visit to pass over at each self-crossing. Reverse
the component order and all traversals to obtain the complementary
request at every crossing. One request therefore costs at most
$\lfloor c/2\rfloor$ switches. Construct the descending lift of each
component within its own height slab by an increasing affine rescaling
of the above heights, with earlier components occupying higher slabs.
All mixed crossing bits are then correct, and all unknotting triangles
stay inside the slab of the component being simplified. Their terminal
triangles bound disjoint filled disks in those slabs. The height-change
lemma applies to the whole link since its argument is local to a vertex
and to the other branch at a crossing.
\end{proof}

The arithmetic assertions use bounded scans over gaps and visits. The
strip verification uses a loop on the remaining suffix and rational
triangle-intersection tests. In a Lean development, first prove a
one-gap update lemma on a cyclic list in which each label occurs once
over and once under. Next formalize a single-vertex height move and
the two triangles of one descending slab. Induction on the suffix emits
the explicit spatial move list. This result supplies crossing changes
followed by three-dimensional PL moves. A conversion of arbitrary
spatial triangle moves to a movie for the original projection is a
separate theorem. For the classical warping-degree formulation see
A.~Shimizu, \emph{The warping degree of a knot diagram}
(\href{https://arxiv.org/abs/0809.1334}{arXiv:0809.1334}).
